\clearpage
\pdfbookmark[0]{Theory of the Combination of Observations. Second Part}{theoria-posterior}
\thispagestyle{empty}
\SourceLeaf{032}
\begin{center}
\vspace*{2em}
{\Large THEORY\par}
\bigskip
{\Large\bfseries OF THE COMBINATION OF OBSERVATIONS\par}
\bigskip
{\large SUBJECT TO THE SMALLEST ERRORS\par}
\bigskip
{\large\textsc{Second Part}\par}
\bigskip
BY\par
\bigskip
{\large CARL FRIEDRICH GAUSS\par}
\bigskip
{\small PRESENTED TO THE ROYAL SOCIETY OF SCIENCES IN 1823, ON FEBRUARY 2.\par}
\bigskip
\rule{.22\linewidth}{.3pt}\par
\bigskip
Commentationes societatis regiae scientiarum Gottingensis recentiores.
\mbox{Vol.~\textsc{v}.}
\par\medskip
Göttingen MDCCCXXIII.
\par\bigskip
\rule{.22\linewidth}{.3pt}
\end{center}
\clearpage
\renewcommand{\Article}[1]{\par\medskip\Needspace{4\baselineskip}\pdfbookmark[1]{Art. #1}{theoria-posterior-art-#1}\centerline{\textbf{#1.}}\nobreak\smallskip}
\SourceLeaf{034}
\begin{center}
{\large THEORY\par}
\medskip
{\Large\bfseries OF THE COMBINATION OF OBSERVATIONS\par}
\medskip
{\large SUBJECT TO THE SMALLEST ERRORS.\par}
\medskip
{\large\textsc{Second Part.}\par}
\medskip
\rule{.22\linewidth}{.3pt}
\end{center}
\Article{23}
Several investigations still remain by which the preceding theory will be both illustrated and extended.

First of all, we must investigate whether the elimination by means of which the indeterminates \(x,y,z\) etc. are to be expressed through \(\xi,\eta,\zeta\) etc. is always possible. Since the number of the former equals that of the latter, the theory of elimination in linear equations gives the following result about this elimination. If \(\xi,\eta,\zeta\) etc. are mutually independent, it will certainly be possible; otherwise, impossible. Let us suppose for the moment that \(\xi,\eta,\zeta\) etc. are not mutually independent. Suppose instead that an identity exists between them
\[
0=F\xi+G\eta+H\zeta+\text{etc.}+K
\]
We shall therefore have
\[
\begin{aligned}
F\Sigma aa+G\Sigma ab+H\Sigma ac+\text{etc.}&=0\\
F\Sigma ab+G\Sigma bb+H\Sigma bc+\text{etc.}&=0\\
F\Sigma ac+G\Sigma bc+H\Sigma cc+\text{etc.}&=0
\end{aligned}
\]
etc., and also
\[
F\Sigma al+G\Sigma bl+H\Sigma cl+\text{etc.}=-K
\]
\par\SourceLeaf{035}
Putting, further,
\[
\left.\begin{aligned}
aF+bG+cH+\text{etc.}&=\theta\\
a'F+b'G+c'H+\text{etc.}&=\theta'\\
a''F+b''G+c''H+\text{etc.}&=\theta''
\end{aligned}\right\}\quad\text{(I)}
\]
etc., we obtain
\[
\begin{aligned}
a\theta+a'\theta'+a''\theta''+\text{etc.}&=0\\
b\theta+b'\theta'+b''\theta''+\text{etc.}&=0\\
c\theta+c'\theta'+c''\theta''+\text{etc.}&=0
\end{aligned}
\]
etc., and also
\[
l\theta+l'\theta'+l''\theta''+\text{etc.}=-K
\]
Multiplying equations (I) respectively by \(\theta,\theta',\theta''\) etc. and adding, we therefore obtain:
\[
0=\theta\theta+\theta'\theta'+\theta''\theta''+\text{etc.}
\]
This equation evidently cannot hold unless \(\theta=0,\allowbreak\theta'=0,\allowbreak\theta''=0\) etc. hold simultaneously. From this we first conclude that necessarily \(K=0\). Equations (I) then show that the functions \(v,v',v''\) etc. are such that their values do not change if the values of the quantities \(x,y,z\) etc. receive increments or decrements respectively proportional to \(F,G,H\) etc.; evidently the same will hold for the functions \(V,V',V''\) etc. Our supposition can therefore hold only in a case in which even from the exact values of the quantities \(V,V',V''\) etc. it would have been impossible to determine the values of the unknowns \(x,y,z\) etc., that is, in which the problem would by its nature have been indeterminate, a case we excluded from our investigation.
\Article{24}
Let \(\beta,\beta',\beta''\) etc. denote multipliers which have the same relation to the indeterminate \(y\) as \(\alpha,\alpha',\alpha''\) etc. have to \(x\); namely, let
\[
\begin{aligned}
a[\beta\alpha]+b[\beta\beta]+c[\beta\gamma]+\text{etc.}&=\beta\\
a'[\beta\alpha]+b'[\beta\beta]+c'[\beta\gamma]+\text{etc.}&=\beta'\\
a''[\beta\alpha]+b''[\beta\beta]+c''[\beta\gamma]+\text{etc.}&=\beta''
\end{aligned}
\]
etc., so that identically
\[
\beta v+\beta'v'+\beta''v''+\text{etc.}=y-B
\]
\par\SourceLeaf{036}
Similarly, let \(\gamma,\gamma',\gamma''\) etc. be analogous multipliers with respect to the indeterminate \(z\), namely
\[
\begin{aligned}
a[\gamma\alpha]+b[\gamma\beta]+c[\gamma\gamma]+\text{etc.}&=\gamma\\
a'[\gamma\alpha]+b'[\gamma\beta]+c'[\gamma\gamma]+\text{etc.}&=\gamma'\\
a''[\gamma\alpha]+b''[\gamma\beta]+c''[\gamma\gamma]+\text{etc.}&=\gamma''
\end{aligned}
\]
etc., so that identically
\[
\gamma v+\gamma'v'+\gamma''v''+\text{etc.}=z-C
\]
and so on. In this way, just as we already found in art. 20 that
\[
\Sigma\alpha a=1,\quad\Sigma\alpha b=0,\quad\Sigma\alpha c=0,\quad\text{etc., and also}\quad\Sigma\alpha l=-A
\]
we shall also have
\[
\begin{aligned}
\Sigma\beta a&=0,\quad\Sigma\beta b=1,\quad\Sigma\beta c=0,\quad\text{etc., and}\quad\Sigma\beta l=-B\\
\Sigma\gamma a&=0,\quad\Sigma\gamma b=0,\quad\Sigma\gamma c=1,\quad\text{etc., and}\quad\Sigma\gamma l=-C
\end{aligned}
\]
and so on. Likewise, just as art. 20 gave \(\Sigma\alpha\alpha=[\alpha\alpha]\), we shall also have
\[
\Sigma\beta\beta=[\beta\beta],\quad\Sigma\gamma\gamma=[\gamma\gamma]\quad\text{etc.}
\]
Multiplying, further, the values of \(\alpha,\alpha',\alpha''\) etc. (art. 20. IV) respectively by \(\beta,\beta',\beta''\) etc. and adding, we obtain
\[
\alpha\beta+\alpha'\beta'+\alpha''\beta''\quad\text{etc.}=[\alpha\beta],\quad\text{or}\quad\Sigma\alpha\beta=[\alpha\beta]
\]
But multiplying the values of \(\beta,\beta',\beta''\) etc. respectively by \(\alpha,\alpha',\alpha''\) etc. and adding gives, in the same way,
\[
\alpha\beta+\alpha'\beta'+\alpha''\beta''+\text{etc.}=[\beta\alpha],\quad\text{hence}\quad[\alpha\beta]=[\beta\alpha]
\]
In precisely the same way we obtain
\[
[\alpha\gamma]=[\gamma\alpha]=\Sigma\alpha\gamma,\quad[\beta\gamma]=[\gamma\beta]=\Sigma\beta\gamma\quad\text{etc.}
\]
\clearpage
